\documentclass[11pt,a4paper]{article}
\usepackage[utf8]{inputenc}
\usepackage[vietnamese]{babel}
\usepackage{amsmath,amssymb}
\usepackage[version=4]{mhchem}
\usepackage{geometry}
\geometry{top=1.5cm,bottom=1.2cm,left=1.4cm,right=1.4cm,headheight=28pt,headsep=12pt,footskip=18pt}
\usepackage{graphicx}
\usepackage{xcolor}
\usepackage{tabularx}
\usepackage{array}
\usepackage{tikz}
\usepackage{fancyhdr}
\usepackage{enumitem}
\usepackage{colortbl}
\usepackage{tcolorbox}
\usepackage{multido}

\definecolor{hfRed}{HTML}{D32F2F}
\definecolor{hfGreen}{HTML}{8BC390}
\definecolor{mainGreen}{HTML}{2E7D32}
\definecolor{solBlue}{HTML}{1565C0}
\definecolor{tableGreen}{HTML}{E8F5E9}

\usepackage{ex_test}

\pagestyle{fancy}
\fancyhf{}
\renewcommand{\headrulewidth}{0.75pt}
\renewcommand{\headrule}{\hbox to\headwidth{\color{hfGreen}\leaders\hrule height \headrulewidth\hfill}}
\renewcommand{\footrulewidth}{0.75pt}
\renewcommand{\footrule}{\hbox to\headwidth{\color{hfGreen}\leaders\hrule height \footrulewidth\hfill}}

\fancyhead[L]{\small\itshape\color{hfRed}Tài liệu lưu hành nội bộ}
\fancyhead[R]{\small\color{hfRed}Thầy \textbf{TRẦN VĂN THẠNH} - 0777.470.803}

\fancyfoot[L]{\footnotesize\itshape\color{hfRed}CS1: 6/15 Nguyễn Hoàng, P. Kim Long, TP Huế \quad|\quad CS2: 24 Đặng Thái Thân, TP Huế}
\fancyfoot[R]{\small\textbf{\color{hfRed}Trang \thepage}}

\setlength{\parindent}{0pt}
\setlength{\parskip}{2pt}

\begin{document}

\begin{minipage}[t]{0.48\textwidth}
\textbf{SỞ GD\&ĐT THÀNH PHỐ HUẾ}\\
\textbf{TRƯỜNG THPT BÙI THỊ XUÂN}\\
\textit{ĐỀ CHÍNH THỨC}\\
\textit{(Đề thi có 02 trang)}
\end{minipage}
\hfill
\begin{minipage}[t]{0.48\textwidth}
\raggedleft
\textbf{\color{hfRed}KIỂM TRA ĐÁNH GIÁ GIỮA KÌ}\\
\textbf{HỌC KÌ I, NĂM HỌC 2024 - 2025}\\
\textbf{MÔN: HOÁ HỌC\_LỚP 10}\\
\textit{Thời gian làm bài: 45 phút (không kể thời gian phát đề)}\\
\textbf{MÃ ĐỀ: 1004}
\end{minipage}

\vspace{4pt}
\noindent
\textit{Họ, tên thí sinh:} \dotfill \textit{Lớp:} \dotfill \textit{Số báo danh:} \dotfill

\vspace{3pt}
\hrule color{hfGreen} height 0.5pt

\vspace{4pt}
\noindent
{\color{mainGreen}\textbf{\large PHẦN I. Câu trắc nghiệm nhiều lựa chọn (3,0 điểm)}}\\
\textit{\small Học sinh trả lời từ câu 1 đến câu 12. Mỗi câu hỏi học sinh chỉ chọn một phương án.}
\vspace{2pt}

\begin{ex}
Phát biểu nào sau đây không đúng?
\choice
{Nguyên tử được cấu thành từ các hạt cơ bản là proton, neutron và electron}
{Hạt nhân nguyên tử cấu thành từ các hạt proton và neutron}
{\True Nguyên tử có cấu trúc đặc khít, gồm vỏ nguyên tử và hạt nhân nguyên tử}
{Vỏ nguyên tử cấu thành từ các hạt electron}
\end{ex}
\begin{ex}
Hạt có khối lượng xấp xỉ 1 amu trong hạt nhân nguyên tử là
\choice
{proton và electron}
{neutron và electron}
{electron}
{\True proton và neutron}
\end{ex}
\begin{ex}
Đường kính nguyên tử khoảng $10^{-1}\text{ nm}$, đường kính hạt nhân khoảng $10^{-5}\text{ nm}$. Đường kính nguyên tử gấp khoảng bao nhiêu lần đường kính hạt nhân?
\choice
{10 lần}
{100 lần}
{1 000 lần}
{\True 10 000 lần}
\end{ex}
\begin{ex}
Hoạt động nào sau đây thuộc phương pháp thực nghiệm trong học tập và nghiên cứu hóa học?
\choice
{Đọc tài liệu và ghi nhớ nội dung lí thuyết}
{\True Tiến hành thí nghiệm, quan sát hiện tượng và rút ra kết luận}
{Xem video bài giảng của giáo viên}
{Trao đổi nhóm để tóm tắt nội dung bài học}
\end{ex}
\begin{ex}
Để kiểm tra một giả thuyết (dự đoán) trong nghiên cứu hóa học, người ta thường
\choice
{chỉ dựa vào ý kiến của số đông}
{chỉ đọc lại kết luận có sẵn trong sách}
{\True tiến hành thí nghiệm, quan sát hiện tượng, thu thập số liệu rồi đối chiếu với dự đoán}
{chọn những kết quả phù hợp với dự đoán và bỏ qua những kết quả khác}
\end{ex}
\begin{ex}
Cho ba nguyên tử: $^{12}_6\ce{X}$, $^{14}_6\ce{Y}$ và $^{14}_7\ce{Z}$. Nhận định nào sau đây đúng?
\choice
{Y và Z là hai đồng vị của cùng một nguyên tố vì có cùng số khối}
{\True X và Y là hai đồng vị của cùng một nguyên tố hóa học}
{X và Z có cùng số neutron}
{X và Z có cùng số khối}
\end{ex}
\begin{ex}
Mô hình nguyên tử hiện đại mô tả chuyển động của electron là
\choice
{theo quỹ đạo hình tròn hay bầu dục}
{theo quỹ đạo đường thẳng}
{\True không theo quỹ đạo xác định}
{theo quỹ đạo hình cầu}
\end{ex}
\begin{ex}
Nhà khoa học nào sau đây đã công bố bảng tuần hoàn các nguyên tố hóa học đầu tiên vào năm 1869?
\choice
{\True D. I. Mendeleev}
{E. Rutherford}
{J. J. Thomson}
{N. Bohr}
\end{ex}
\begin{ex}
Nguyên tử của nguyên tố X có tổng số electron trong các phân lớp p là 10. Vị trí của X trong bảng tuần hoàn là
\choice
{ô số 10, chu kì 2, nhóm VIIIA}
{\True ô số 16, chu kì 3, nhóm VIA}
{ô số 14, chu kì 3, nhóm IVA}
{ô số 16, chu kì 3, nhóm IVA}
\end{ex}
\begin{ex}
Cho các nguyên tố $\ce{X}\ (Z = 13)$, $\ce{Y}\ (Z = 17)$ và $\ce{T}\ (Z = 20)$. Phát biểu nào sau đây đúng?
\choice
{X, Y đều là phi kim}
{\True X, T là kim loại; Y là phi kim}
{X là phi kim; Y, T là kim loại}
{X, Y, T đều là kim loại}
\end{ex}
\begin{ex}
Số hiệu nguyên tử của nguyên tố hóa học bằng
\choice
{\True số thứ tự của ô nguyên tố}
{số thứ tự của chu kì}
{số thứ tự của nhóm}
{số electron lớp ngoài cùng của nguyên tử}
\end{ex}
\begin{ex}
Nhóm A bao gồm các nguyên tố
\choice
{s}
{p}
{d và f}
{\True s và p}
\end{ex}
\vspace{4pt}
\noindent
{\color{mainGreen}\textbf{\large PHẦN II. Câu trắc nghiệm đúng – sai (2,0 điểm)}}\\
\textit{\small Học sinh trả lời câu 1 và câu 2. Trong mỗi ý a), b), c), d) ở mỗi câu, học sinh chọn đúng hoặc sai.}
\vspace{2pt}

\begin{ex}
Magnesium được dùng để chế tạo hợp kim nhẹ, calcium được dùng trong sản xuất vôi và xi măng. Cho nguyên tử magnesium ($\ce{Mg}, Z = 12$) và calcium ($\ce{Ca}, Z = 20$).
\choiceTF
{Mg là nguyên tố s, còn Ca là nguyên tố p}
{\True Từ cấu hình electron, xác định được Ca nằm ở ô số 20, chu kì 4, nhóm IIA}
{\True Bảng tuần hoàn hiện nay có 8 nhóm A và 8 nhóm B}
{Mg và Ca thuộc cùng một chu kì vì nguyên tử của chúng đều có 2 electron ở lớp ngoài cùng}
\end{ex}
\begin{ex}
Nguyên tử của nguyên tố X có cấu hình electron kết thúc bằng $3p^2$. Nguyên tử của nguyên tố Y có cấu hình electron kết thúc bằng $3p^5$.
\choiceTF
{\True Nguyên tử Y có 1 electron độc thân ở trạng thái cơ bản}
{\True Nguyên tử Y có 7 electron ở lớp ngoài cùng nên Y là phi kim}
{Phân lớp 3p có 3 orbital và chứa tối đa 3 electron}
{Cấu hình electron theo lớp của X là (2,8,2)}
\end{ex}
\vspace{4pt}
\noindent
{\color{mainGreen}\textbf{\large PHẦN III. Câu trắc nghiệm trả lời ngắn (2,0 điểm)}}\\
\textit{\small Học sinh trả lời từ câu 1 đến câu 4.}
\vspace{2pt}

\begin{ex}
Cho các nguyên tố có số hiệu nguyên tử lần lượt là 8, 12, 14, 16 và 19. Trong số đó có hai nguyên tố thuộc cùng một nhóm. Tổng số hiệu nguyên tử của hai nguyên tố này là bao nhiêu?

\noindent\textit{Trả lời:} \dotfill
\end{ex}
\begin{ex}
Nguyên tử Y có lớp electron thứ nhất đã bão hòa và lớp electron thứ hai chứa 5 electron. Số electron độc thân của Y là bao nhiêu?

\noindent\textit{Trả lời:} \dotfill
\end{ex}
\begin{ex}
Nguyên tử R có điện tích lớp vỏ là $-33{,}642\cdot 10^{-19}\text{ C}$. Biết điện tích của một electron là $-1{,}602\cdot 10^{-19}\text{ C}$. Nguyên tử R có bao nhiêu electron?

\noindent\textit{Trả lời:} \dotfill
\end{ex}
\begin{ex}
Trong tự nhiên, đồng có hai đồng vị bền là $^{63}\ce{Cu}$ và $^{65}\ce{Cu}$, nguyên tử khối trung bình của đồng là 63,54. Phần trăm số nguyên tử của đồng vị $^{65}\ce{Cu}$ là bao nhiêu?

\noindent\textit{Trả lời:} \dotfill
\end{ex}
\vspace{6pt}
\noindent
{\color{mainGreen}\textbf{\large PHẦN IV. Tự luận (3,0 điểm)}}\\
\textit{\small Học sinh trình bày lời giải chi tiết cho các câu hỏi sau.}
\vspace{2pt}

\noindent
\textbf{Câu 1. (1,0 điểm)} Nguyên tử của nguyên tố Y có tổng số hạt là 10 và số khối bằng 7.\\
a) Xác định số proton, số neutron, số electron của Y và viết kí hiệu nguyên tử của Y.\\
b) Cho thêm hai nguyên tố T và V có kí hiệu lần lượt là $^{6}_{3}\ce{T}$ và $^{7}_{4}\ce{V}$. Trong ba nguyên tử Y, T, V, những nguyên tử nào là đồng vị của cùng một nguyên tố? Giải thích.

\vspace{1pt}
\noindent
\multido{\n=1+1}{7}{\dotfill\\[4pt]}

\vspace{4pt}
\noindent
\textbf{Câu 2. (1,0 điểm)} Nguyên tử của nguyên tố X có 3 lớp electron và lớp thứ ba chứa 6 electron.\\
a) Xác định số hiệu nguyên tử của X và viết cấu hình electron nguyên tử của X (theo lớp và phân lớp).\\
b) Biểu diễn sự phân bố electron của X theo ô orbital và cho biết X có bao nhiêu electron độc thân. Nếu nguyên tử X nhận thêm electron để đạt cấu hình bền thì cấu hình electron của ion tạo thành giống cấu hình của nguyên tử khí hiếm nào?

\vspace{1pt}
\noindent
\multido{\n=1+1}{8}{\dotfill\\[4pt]}

\vspace{4pt}
\noindent
\textbf{Câu 3. (1,0 điểm)} X và Y là hai nguyên tố thuộc cùng một chu kì và thuộc hai nhóm A liên tiếp trong bảng tuần hoàn, $Z_{\ce{X}} < Z_{\ce{Y}}$. Tổng số hiệu nguyên tử của X và Y bằng 25.\\
a) Hãy xác định số hiệu nguyên tử của X và Y.\\
b) Viết cấu hình electron của X, Y và xác định vị trí của mỗi nguyên tố trong bảng tuần hoàn.\\
c) Dựa vào cấu hình electron lớp ngoài cùng, cho biết X, Y là kim loại hay phi kim và thuộc loại nguyên tố s, p, d hay f.

\vspace{1pt}
\noindent
\multido{\n=1+1}{9}{\dotfill\\[4pt]}

\vspace{6pt}
\begin{center}
\textbf{\large ---------- HẾT ----------}
\end{center}

\end{document}
